% AMC 12A 2023 Problem 18 Strategic Analysis
% Generated using Strategic Problem Solving Framework v2.1

\documentclass[11pt]{article}
\usepackage{amsmath, amssymb, amsthm}
\usepackage{geometry}
\usepackage{enumitem}
\usepackage{tcolorbox}
\tcbuselibrary{skins,breakable}
\usepackage{xcolor}

% Page setup
\geometry{margin=1in}
\setlength{\parindent}{0pt}
\setlength{\parskip}{6pt}

% Custom colored boxes
\newtcolorbox{problembox}[1][]{
    colback=blue!5!white,
    colframe=blue!75!black,
    title=Problem Configuration Analysis,
    breakable,
    #1
}

\newtcolorbox{strategybox}[1][]{
    colback=pink!5!white,
    colframe=pink!75!black,
    title=Strategic Framework Application,
    breakable,
    #1
}

\newtcolorbox{tacticalbox}[1][]{
    colback=green!5!white,
    colframe=green!75!black,
    title=Tactical Selection,
    breakable,
    #1
}

\newtcolorbox{conversionbox}[1][]{
    colback=yellow!5!white,
    colframe=yellow!75!black,
    title=Conversion Techniques,
    breakable,
    #1
}

\newtcolorbox{verificationbox}[1][]{
    colback=orange!5!white,
    colframe=orange!75!black,
    title=Verification Strategy,
    breakable,
    #1
}

\newtcolorbox{roadmapbox}[1][]{
    colback=purple!5!white,
    colframe=purple!75!black,
    title=Solution Roadmap,
    breakable,
    #1
}

\newtcolorbox{competitionbox}[1][]{
    colback=red!5!white,
    colframe=red!75!black,
    title=Competition Strategy,
    breakable,
    #1
}

\newtcolorbox{pitfallbox}[1][]{
    colback=gray!5!white,
    colframe=gray!75!black,
    title=Common Pitfalls and Warnings,
    breakable,
    #1
}

\newtcolorbox{solutionbox}[1][]{
    colback=cyan!5!white,
    colframe=cyan!75!black,
    title=Detailed Solution,
    breakable,
    #1
}

\newtcolorbox{keyinsightbox}[1][]{
    colback=magenta!5!white,
    colframe=magenta!75!black,
    title=Key Insights,
    breakable,
    #1
}

\begin{document}

\title{\textbf{AMC 12A 2023 Problem 18}\\Strategic Analysis}
\author{Using Framework v2.1}
\date{}
\maketitle

\section*{Problem Statement}

Circle $C_1$ and $C_2$ each have radius $1$, and the distance between their centers is $\frac{1}{2}$. Circle $C_3$ is the largest circle internally tangent to both $C_1$ and $C_2$. Circle $C_4$ is internally tangent to both $C_1$ and $C_2$ and externally tangent to $C_3$. What is the radius of $C_4$?

\[\textbf{(A) } \frac{1}{14} \qquad \textbf{(B) } \frac{1}{12} \qquad \textbf{(C) } \frac{1}{10} \qquad \textbf{(D) } \frac{3}{28} \qquad \textbf{(E) } \frac{1}{9}\]

\begin{problembox}
\section*{1. Problem Configuration Analysis}

\subsection*{A. Problem Classification}
\begin{itemize}[leftmargin=*]
    \item \textbf{Competition}: AMC 12A 2023
    \item \textbf{Problem Number}: 18
    \item \textbf{Difficulty Band}: Upper-mid to late-mid range (Problem 18/25)
    \item \textbf{Expected Time}: 4-6 minutes
    \item \textbf{Point Value}: 6 points (standard AMC 12 scoring)
\end{itemize}

\subsection*{B. Mathematical Domain Classification}
\begin{itemize}[leftmargin=*]
    \item \textbf{Primary Domain}: Geometry (Circles)
    \item \textbf{Secondary Domain}: Coordinate Geometry / Tangency
    \item \textbf{Tertiary Domain}: Pythagorean Theorem
    \item \textbf{Cross-Domain Potential}: Moderate (geometry, algebra, spatial reasoning)
\end{itemize}

\subsection*{C. Problem Type}
\begin{itemize}[leftmargin=*]
    \item \textbf{Format}: Multiple choice (numerical answer)
    \item \textbf{Question Type}: ``What is the radius...''
    \item \textbf{Core Task}: Find radius of circle with multiple tangency conditions
    \item \textbf{Answer Form}: Simple fraction
\end{itemize}

\subsection*{D. Given Information}
\begin{itemize}[leftmargin=*]
    \item $C_1$ and $C_2$: radius = 1 each
    \item Distance between centers of $C_1$ and $C_2$: $\frac{1}{2}$
    \item $C_3$: largest circle internally tangent to both $C_1$ and $C_2$
    \item $C_4$: internally tangent to both $C_1$ and $C_2$, externally tangent to $C_3$
    \item Two large circles overlap significantly (centers $\frac{1}{2}$ apart, radii 1 each)
\end{itemize}

\subsection*{E. Constraints}
\begin{itemize}[leftmargin=*]
    \item $C_3$ must be the \textit{largest} such circle (unique)
    \item $C_4$ satisfies three tangency conditions simultaneously
    \item All circles lie in the same plane
\end{itemize}

\subsection*{F. Objective}
\begin{itemize}[leftmargin=*]
    \item \textbf{Find}: Radius of circle $C_4$
    \item \textbf{Output}: Select correct fraction from choices
\end{itemize}

\subsection*{G. Key Structural Features}
\begin{itemize}[leftmargin=*]
    \item \textbf{Symmetry}: Configuration has reflection symmetry about perpendicular bisector
    \item \textbf{Tangency conditions}: Internal tangency means distance between centers = $|r_1 - r_2|$
    \item \textbf{External tangency}: Distance between centers = $r_1 + r_2$
    \item \textbf{Coordinate setup}: Natural to place centers of $C_1$ and $C_2$ symmetrically
    \item \textbf{Right triangle}: Perpendicular bisector creates right triangle
\end{itemize}

\subsection*{H. Strategic Orientation}
\begin{itemize}[leftmargin=*]
    \item \textbf{Primary Strategy}: Set up coordinate system, use tangency conditions
    \item \textbf{Key Technique}: Pythagorean theorem on right triangle
    \item \textbf{Expected Difficulty}: Moderate to challenging - requires careful setup
    \item \textbf{Time Pressure}: Should be completed in 4-6 minutes
\end{itemize}
\end{problembox}

\begin{strategybox}
\section*{2. Strategic Framework Application}

\subsection*{A. Psychological Strategies}
\begin{itemize}[leftmargin=*]
    \item \textbf{Confidence Calibration}: 
    \begin{itemize}
        \item Upper-mid difficulty problem
        \item Standard geometry techniques apply
        \item Answer choices are simple fractions
    \end{itemize}
    \item \textbf{Pattern Recognition}:
    \begin{itemize}
        \item Recognize tangency conditions translate to distance formulas
        \item Symmetry suggests coordinate placement
        \item Multiple tangencies create system of constraints
    \end{itemize}
    \item \textbf{Problem Familiarity}:
    \begin{itemize}
        \item Classic tangent circles problem
        \item Similar to Apollonian gasket configurations
        \item Standard competition geometry
    \end{itemize}
\end{itemize}

\subsection*{B. Investigation Strategies}
\begin{itemize}[leftmargin=*]
    \item \textbf{Initial Exploration}:
    \begin{itemize}
        \item Draw diagram carefully (provided in problem)
        \item Identify centers and radii relationships
        \item Note that $C_1$ and $C_2$ overlap significantly
    \end{itemize}
    \item \textbf{Coordinate Setup}:
    \begin{itemize}
        \item Place origin at midpoint between centers of $C_1$ and $C_2$
        \item $C_1$ centered at $(-\frac{1}{4}, 0)$, $C_2$ centered at $(\frac{1}{4}, 0)$
        \item By symmetry, $C_3$ and $C_4$ centered on $y$-axis
    \end{itemize}
    \item \textbf{Key Calculations}:
    \begin{itemize}
        \item Find radius of $C_3$ first
        \item Set up equation for radius of $C_4$ using tangency conditions
        \item Use Pythagorean theorem
    \end{itemize}
\end{itemize}

\subsection*{C. Argument Construction}
\begin{itemize}[leftmargin=*]
    \item \textbf{Logical Structure}:
    \begin{enumerate}
        \item Set up coordinates with origin at midpoint
        \item Find radius of $C_3$ using internal tangency to $C_1$ and $C_2$
        \item Let radius of $C_4$ be $r$, find its center location
        \item Use internal tangency to $C_1$ (or $C_2$): distance equation
        \item Use external tangency to $C_3$: distance equation
        \item Apply Pythagorean theorem to right triangle
        \item Solve for $r$
    \end{enumerate}
    \item \textbf{Key Deductions}:
    \begin{itemize}
        \item Radius of $C_3$ is $\frac{3}{4}$ (from internal tangency conditions)
        \item Center of $C_4$ lies on positive $y$-axis
        \item Distance from origin to boundary of $C_3$ is $\frac{3}{4}$
    \end{itemize}
\end{itemize}

\subsection*{D. Cross-Domain Translation}
\begin{itemize}[leftmargin=*]
    \item \textbf{Geometric}: Tangent circles and distance relationships
    \item \textbf{Algebraic}: System of distance equations
    \item \textbf{Coordinate}: Place circles in coordinate system
    \item \textbf{Trigonometric}: Right triangle with Pythagorean theorem
\end{itemize}
\end{strategybox}

\begin{tacticalbox}
\section*{3. Tactical Methodology Selection}

\subsection*{A. Core Mathematical Tactics}
\begin{itemize}[leftmargin=*]
    \item \textbf{Coordinate Geometry}: Place centers strategically
    \item \textbf{Tangency Conditions}: Translate to distance formulas
    \item \textbf{Symmetry}: Use reflection symmetry to simplify
    \item \textbf{Pythagorean Theorem}: Apply to right triangle
    \item \textbf{Algebraic Manipulation}: Solve for unknown radius
\end{itemize}

\subsection*{B. Domain-Specific Tactics}
\begin{itemize}[leftmargin=*]
    \item \textbf{Circle Tangency}:
    \begin{itemize}
        \item Internal tangency: $d = |r_1 - r_2|$ (larger circle contains smaller)
        \item External tangency: $d = r_1 + r_2$
        \item Distance between centers determines relationship
    \end{itemize}
    \item \textbf{Coordinate Setup}:
    \begin{itemize}
        \item Origin at midpoint: centers at $(\pm\frac{1}{4}, 0)$
        \item Symmetry implies $C_3$ and $C_4$ on $y$-axis
        \item Right triangle naturally appears
    \end{itemize}
    \item \textbf{Problem-Solving Sequence}:
    \begin{itemize}
        \item Find $C_3$ first (largest, well-defined)
        \item Use $C_3$ to constrain $C_4$
        \item Three tangency conditions → solvable system
    \end{itemize}
\end{itemize}

\subsection*{C. Tactical Priorities}
\begin{enumerate}[leftmargin=*]
    \item Set up coordinate system (60 seconds)
    \item Find radius of $C_3$ (60 seconds)
    \item Set up equation for $C_4$ using tangency (90 seconds)
    \item Apply Pythagorean theorem (60 seconds)
    \item Solve for $r$ (60 seconds)
    \item Verify answer (30 seconds)
\end{enumerate}
\end{tacticalbox}

\begin{conversionbox}
\section*{4. Conversion Techniques Application}

\subsection*{A. Method 1: Coordinate Geometry (Standard Approach)}

\textbf{Step 1: Set up coordinates}
\begin{itemize}
    \item Place origin $O$ at midpoint between centers of $C_1$ and $C_2$
    \item Center of $C_1$: $A = (-\frac{1}{4}, 0)$
    \item Center of $C_2$: $B = (\frac{1}{4}, 0)$
    \item Both have radius 1
\end{itemize}

\textbf{Step 2: Find radius of $C_3$}
\begin{itemize}
    \item $C_3$ is internally tangent to both $C_1$ and $C_2$
    \item By symmetry, center of $C_3$ is at origin $(0, 0)$
    \item Let radius of $C_3$ be $r_3$
    \item Internal tangency to $C_1$: distance from $O$ to $A$ equals $1 - r_3$
    \[\frac{1}{4} = 1 - r_3 \implies r_3 = \frac{3}{4}\]
\end{itemize}

\textbf{Step 3: Find location of center of $C_4$}
\begin{itemize}
    \item By symmetry, center of $C_4$ is on $y$-axis at $(0, y_0)$ for some $y_0 > 0$
    \item Let radius of $C_4$ be $r$
\end{itemize}

\textbf{Step 4: Apply tangency conditions}
\begin{itemize}
    \item Internal tangency to $C_1$: 
    \[\sqrt{(0-(-\frac{1}{4}))^2 + (y_0-0)^2} = 1 - r\]
    \[\sqrt{\frac{1}{16} + y_0^2} = 1 - r\]
    
    \item External tangency to $C_3$:
    \[y_0 = r_3 + r = \frac{3}{4} + r\]
\end{itemize}

\textbf{Step 5: Substitute and solve}
\begin{itemize}
    \item Substitute $y_0 = \frac{3}{4} + r$:
    \[\sqrt{\frac{1}{16} + (\frac{3}{4} + r)^2} = 1 - r\]
    
    \item Square both sides:
    \[\frac{1}{16} + (\frac{3}{4} + r)^2 = (1-r)^2\]
    
    \item Expand:
    \[\frac{1}{16} + \frac{9}{16} + \frac{3r}{2} + r^2 = 1 - 2r + r^2\]
    
    \item Simplify:
    \[\frac{10}{16} + \frac{3r}{2} = 1 - 2r\]
    \[\frac{5}{8} + \frac{3r}{2} = 1 - 2r\]
    \[\frac{3r}{2} + 2r = 1 - \frac{5}{8} = \frac{3}{8}\]
    \[\frac{7r}{2} = \frac{3}{8}\]
    \[r = \frac{3}{8} \cdot \frac{2}{7} = \frac{3}{28}\]
\end{itemize}

\subsection*{B. Method 2: Direct Right Triangle Approach}

\textbf{Setup}: Same coordinate system as Method 1.

\textbf{Key observation}: If we draw radii from centers $A$, $O$, and center of $C_4$ (call it $D$), we form a right triangle.

\begin{itemize}
    \item $A = (-\frac{1}{4}, 0)$ (center of $C_1$)
    \item $O = (0, 0)$ (center of $C_3$)
    \item $D = (0, \frac{3}{4} + r)$ (center of $C_4$)
    
    \item Triangle $AOD$ has:
    \begin{itemize}
        \item Right angle at $O$
        \item $AO = \frac{1}{4}$ (horizontal leg)
        \item $OD = \frac{3}{4} + r$ (vertical leg)
        \item $AD = 1 - r$ (hypotenuse, from internal tangency to $C_1$)
    \end{itemize}
    
    \item Pythagorean theorem:
    \[\left(\frac{1}{4}\right)^2 + \left(\frac{3}{4} + r\right)^2 = (1-r)^2\]
\end{itemize}

This gives the same equation as Method 1, yielding $r = \frac{3}{28}$.

\subsection*{C. Method 3: Estimation (Quick Check)}

\begin{itemize}
    \item Radius of $C_3$ is $\frac{3}{4}$
    \item $C_4$ is much smaller than $C_3$
    \item Looking at answer choices: (A) $\frac{1}{14} \approx 0.071$, (D) $\frac{3}{28} \approx 0.107$
    \item From diagram, $r_{C_4}$ looks to be roughly $\frac{1}{10}$ of the large circles
    \item $\frac{3}{28} \approx 0.107$ is close to $\frac{1}{10} = 0.1$
    \item This matches (D)
\end{itemize}

\subsection*{D. Answer Selection}
\begin{itemize}
    \item Calculated: $r = \frac{3}{28}$
    \item Matches choice: \textbf{(D)}
\end{itemize}
\end{conversionbox}

\begin{verificationbox}
\section*{5. Verification Strategy Design}

\subsection*{A. Verification Level}
\begin{itemize}[leftmargin=*]
    \item \textbf{Selected Level}: Level 4 - Direct Substitution
    \item \textbf{Reasoning}: Can verify all tangency conditions
    \item \textbf{Time Budget}: 60 seconds
\end{itemize}

\subsection*{B. Verification Checklist}
\begin{itemize}[leftmargin=*]
    \item \textbf{Check arithmetic}:
    \begin{itemize}
        \item $\frac{7r}{2} = \frac{3}{8}$
        \item $r = \frac{3}{28}$ $\checkmark$
    \end{itemize}
    \item \textbf{Verify tangency to $C_1$}:
    \begin{itemize}
        \item Center of $C_4$: $(0, \frac{3}{4} + \frac{3}{28}) = (0, \frac{21+3}{28}) = (0, \frac{24}{28}) = (0, \frac{6}{7})$
        \item Distance to $A = (-\frac{1}{4}, 0)$:
        \[\sqrt{\frac{1}{16} + \frac{36}{49}} = \sqrt{\frac{49 + 576}{784}} = \sqrt{\frac{625}{784}} = \frac{25}{28}\]
        \item Should equal $1 - r = 1 - \frac{3}{28} = \frac{25}{28}$ $\checkmark$
    \end{itemize}
    \item \textbf{Verify tangency to $C_3$}:
    \begin{itemize}
        \item Distance between centers: $\frac{6}{7}$
        \item Should equal $r_3 + r = \frac{3}{4} + \frac{3}{28} = \frac{21+3}{28} = \frac{24}{28} = \frac{6}{7}$ $\checkmark$
    \end{itemize}
\end{itemize}

\subsection*{C. Alternative Verification Methods}
\begin{itemize}[leftmargin=*]
    \item \textbf{Method 1}: Check Pythagorean theorem directly
    \item \textbf{Method 2}: Verify answer is among choices
    \item \textbf{Method 3}: Dimensional analysis (radius should be small)
\end{itemize}

\subsection*{D. Answer Choice Validation}
\begin{itemize}[leftmargin=*]
    \item Our answer: $\frac{3}{28}$
    \item Matches choice: \textbf{(D)}
    \item Why not the others?
    \begin{itemize}
        \item (A) $\frac{1}{14} = \frac{2}{28}$: too small
        \item (B) $\frac{1}{12} \approx 0.083$: too small
        \item (C) $\frac{1}{10} = 0.1$: close but not exact
        \item (E) $\frac{1}{9} \approx 0.111$: too large
    \end{itemize}
\end{itemize}
\end{verificationbox}

\begin{roadmapbox}
\section*{6. Solution Roadmap}

\subsection*{A. High-Level Solution Path}
\begin{enumerate}[leftmargin=*]
    \item \textbf{Set up coordinates}: Origin at midpoint
    \item \textbf{Find $r_3$}: Use internal tangency to $C_1$ and $C_2$
    \item \textbf{Locate $C_4$}: By symmetry, on $y$-axis
    \item \textbf{Apply Pythagorean theorem}: Use right triangle
    \item \textbf{Solve for $r$}: Get $\frac{3}{28}$
\end{enumerate}

\subsection*{B. Detailed Solution Steps}

\textbf{Step 1}: Set up coordinate system.
\begin{itemize}
    \item Origin at midpoint between centers
    \item $C_1$ centered at $(-\frac{1}{4}, 0)$, radius 1
    \item $C_2$ centered at $(\frac{1}{4}, 0)$, radius 1
\end{itemize}

\textbf{Step 2}: Find radius of $C_3$.
\begin{itemize}
    \item $C_3$ centered at origin (by symmetry)
    \item Internal tangency: $\frac{1}{4} = 1 - r_3$
    \item Therefore: $r_3 = \frac{3}{4}$
\end{itemize}

\textbf{Step 3}: Set up for $C_4$.
\begin{itemize}
    \item Let radius of $C_4$ be $r$
    \item Center at $(0, y_0)$ where $y_0 = \frac{3}{4} + r$ (external tangency to $C_3$)
\end{itemize}

\textbf{Step 4}: Apply Pythagorean theorem.
\begin{itemize}
    \item Right triangle with legs $\frac{1}{4}$ and $\frac{3}{4} + r$, hypotenuse $1-r$:
    \[\left(\frac{1}{4}\right)^2 + \left(\frac{3}{4} + r\right)^2 = (1-r)^2\]
\end{itemize}

\textbf{Step 5}: Solve for $r$.
\begin{itemize}
    \item Expand and simplify: $\frac{7r}{2} = \frac{3}{8}$
    \item Therefore: $r = \frac{3}{28}$
\end{itemize}

\subsection*{C. Key Decision Points}
\begin{itemize}[leftmargin=*]
    \item \textbf{Decision 1}: Place origin at midpoint (uses symmetry)
    \item \textbf{Decision 2}: Find $C_3$ first (simplifies problem)
    \item \textbf{Decision 3}: Recognize right triangle structure
\end{itemize}

\subsection*{D. Time Management}
\begin{itemize}[leftmargin=*]
    \item Understanding problem: 30 seconds
    \item Setting up coordinates: 60 seconds
    \item Finding $r_3$: 60 seconds
    \item Setting up equation for $r_4$: 90 seconds
    \item Solving: 60 seconds
    \item Verification: 60 seconds
    \item \textbf{Total}: 5 minutes 30 seconds
\end{itemize}
\end{roadmapbox}

\begin{competitionbox}
\section*{7. Competition Strategy Integration}

\subsection*{A. Problem Triage}
\begin{itemize}[leftmargin=*]
    \item \textbf{Difficulty Assessment}: Upper-medium (Problem 18/25)
    \item \textbf{Confidence Level}: Moderate if comfortable with coordinate geometry
    \item \textbf{Decision}: Worth attempting if you have 5+ minutes
\end{itemize}

\subsection*{B. Time Allocation}
\begin{itemize}[leftmargin=*]
    \item \textbf{Target Time}: 5-6 minutes
    \item \textbf{Maximum Time}: 7 minutes
    \item \textbf{Quick Strategy}: If stuck, estimate from diagram
\end{itemize}

\subsection*{C. Solution Format}
\begin{itemize}[leftmargin=*]
    \item Multiple choice: Select \textbf{(D)} $\frac{3}{28}$
    \item Verification: Check one tangency condition
    \item Bubble answer sheet carefully
\end{itemize}

\subsection*{D. Strategic Considerations}
\begin{itemize}[leftmargin=*]
    \item \textbf{Diagram essential}: Understanding geometry from diagram
    \item \textbf{Symmetry shortcut}: Recognizes centers lie on axes
    \item \textbf{Estimation backup}: Can narrow to (C) or (D) from diagram
    \item \textbf{Verification important}: Easy to make algebra errors
\end{itemize}
\end{competitionbox}

\begin{pitfallbox}
\section*{Common Pitfalls and Warnings}

\subsection*{A. Conceptual Errors}
\begin{itemize}[leftmargin=*]
    \item \textbf{Confusing tangency types}: Internal vs external tangency have different distance formulas
    \item \textbf{Missing $C_3$ calculation}: Must find $r_3 = \frac{3}{4}$ first
    \item \textbf{Wrong coordinate setup}: Not using symmetry to simplify
    \item \textbf{Forgetting which circle}: Mixing up $C_1$, $C_2$, $C_3$, $C_4$
\end{itemize}

\subsection*{B. Computational Errors}
\begin{itemize}[leftmargin=*]
    \item \textbf{Expanding $(r + \frac{3}{4})^2$}: Must be careful with algebra
    \item \textbf{Pythagorean theorem}: Easy to mix up which is hypotenuse
    \item \textbf{Solving linear equation}: $\frac{7r}{2} = \frac{3}{8}$ gives $r = \frac{3}{28}$
    \item \textbf{Simplification}: $\frac{3}{4} + \frac{3}{28} = \frac{21+3}{28} = \frac{24}{28} = \frac{6}{7}$
\end{itemize}

\subsection*{C. Strategic Errors}
\begin{itemize}[leftmargin=*]
    \item \textbf{Not drawing diagram}: Difficult to visualize without picture
    \item \textbf{Overcomplicating}: Using general formulas instead of coordinates
    \item \textbf{Time management}: This problem takes 5-6 minutes minimum
\end{itemize}

\subsection*{D. Warning Signs}
\begin{itemize}[leftmargin=*]
    \item If you get $r > \frac{1}{4}$, something is wrong ($C_4$ is small)
    \item If answer isn't among choices, recheck algebra
    \item If verification fails, revisit tangency conditions
\end{itemize}
\end{pitfallbox}

\begin{solutionbox}
\subsection*{A. Complete Solution}

\textbf{Setup}: Place origin $O$ at the midpoint between the centers of $C_1$ and $C_2$.
\begin{itemize}
    \item Center of $C_1$: $A = (-\frac{1}{4}, 0)$, radius 1
    \item Center of $C_2$: $B = (\frac{1}{4}, 0)$, radius 1
\end{itemize}

\textbf{Find radius of $C_3$}:
\begin{itemize}
    \item By symmetry, $C_3$ is centered at origin
    \item $C_3$ is internally tangent to $C_1$
    \item Distance from $O$ to $A$ is $\frac{1}{4}$
    \item For internal tangency: $\frac{1}{4} = 1 - r_3$
    \item Therefore: $r_3 = \frac{3}{4}$
\end{itemize}

\textbf{Find radius of $C_4$}:
\begin{itemize}
    \item Let radius of $C_4$ be $r$
    \item By symmetry, center of $C_4$ is on $y$-axis
    \item $C_4$ is externally tangent to $C_3$, so center is at $(0, \frac{3}{4} + r)$
    
    \item Consider right triangle $AOD$ where $D$ is center of $C_4$:
    \begin{itemize}
        \item Horizontal leg $AO = \frac{1}{4}$
        \item Vertical leg $OD = \frac{3}{4} + r$
        \item Hypotenuse $AD = 1 - r$ (from internal tangency to $C_1$)
    \end{itemize}
    
    \item Pythagorean theorem:
    \[\left(\frac{1}{4}\right)^2 + \left(\frac{3}{4} + r\right)^2 = (1-r)^2\]
    
    \item Expand:
    \[\frac{1}{16} + \frac{9}{16} + \frac{3r}{2} + r^2 = 1 - 2r + r^2\]
    
    \item Simplify:
    \[\frac{10}{16} + \frac{3r}{2} = 1 - 2r\]
    \[\frac{5}{8} + \frac{3r}{2} + 2r = 1\]
    \[\frac{7r}{2} = \frac{3}{8}\]
    \[r = \frac{3}{28}\]
\end{itemize}

\subsection*{B. Verification}
\begin{itemize}
    \item Center of $C_4$: $(0, \frac{3}{4} + \frac{3}{28}) = (0, \frac{6}{7})$
    \item Distance to $A = (-\frac{1}{4}, 0)$:
    \[\sqrt{\frac{1}{16} + \frac{36}{49}} = \frac{25}{28}\]
    \item Should equal $1 - \frac{3}{28} = \frac{25}{28}$ $\checkmark$
\end{itemize}
\end{solutionbox}

\begin{keyinsightbox}
\textbf{Key Insight}: The key to this problem is recognizing the geometric structure and using symmetry effectively.

\textbf{Critical observations}:
\begin{itemize}
    \item The configuration has bilateral symmetry about the perpendicular bisector
    \item This means $C_3$ and $C_4$ both have centers on the $y$-axis
    \item The problem naturally sets up a right triangle
    \item Finding $C_3$ first ($r_3 = \frac{3}{4}$) is essential for solving for $C_4$
\end{itemize}

\textbf{Tangency conditions translate to distances}:
\begin{itemize}
    \item Internal tangency: If circle with radius $R$ contains circle with radius $r$, and they're internally tangent, then the distance between centers is $R - r$
    \item External tangency: If two circles with radii $r_1$ and $r_2$ are externally tangent, distance between centers is $r_1 + r_2$
\end{itemize}

\textbf{Right triangle structure}:
\begin{itemize}
    \item The perpendicular from the $y$-axis to the center of $C_1$ creates a right angle
    \item This allows us to use the Pythagorean theorem
    \item The three tangency conditions reduce to one equation in one unknown
\end{itemize}

\textbf{Why this approach works}:
\begin{itemize}
    \item Symmetry reduces variables (centers must lie on axes)
    \item Coordinate geometry makes distance calculations straightforward
    \item Pythagorean theorem is simpler than general distance formulas
    \item Finding $C_3$ first provides a constraint for $C_4$
\end{itemize}

\textbf{General principle}: For multiple tangent circle problems:
\begin{itemize}
    \item Use symmetry to reduce complexity
    \item Set up coordinates to align with symmetry axes
    \item Find simpler circles first, then use them to constrain more complex ones
    \item Translate tangency conditions to algebraic distance equations
\end{itemize}
\end{keyinsightbox}

\section*{Final Answer}

The radius of circle $C_4$ is:

\[\boxed{\textbf{(D) } \frac{3}{28}}\]

\textbf{Solution Summary}:
\begin{itemize}
    \item Place origin at midpoint: centers of $C_1$ and $C_2$ at $(\pm\frac{1}{4}, 0)$
    \item Find $C_3$: radius $= \frac{3}{4}$, centered at origin
    \item Set up for $C_4$: radius $r$, centered at $(0, \frac{3}{4} + r)$
    \item Pythagorean theorem: $(\frac{1}{4})^2 + (\frac{3}{4} + r)^2 = (1-r)^2$
    \item Solve: $r = \frac{3}{28}$
\end{itemize}

\section*{Confidence Assessment}
\begin{itemize}[leftmargin=*]
    \item \textbf{Confidence Level}: 95\%
    \begin{itemize}
        \item Standard coordinate geometry approach
        \item All tangency conditions verified
        \item Arithmetic checked carefully
        \item Answer matches expected size
    \end{itemize}
    \item \textbf{Verification Applied}: Level 4 - Direct Substitution
    \begin{itemize}
        \item Verified internal tangency to $C_1$
        \item Verified external tangency to $C_3$
        \item Checked arithmetic multiple times
    \end{itemize}
    \item \textbf{Time-Confidence Trade-off}: High confidence in 5-6 minutes
\end{itemize}

\end{document}

